% Contextual printed508: blank reverse; principal515, Google519.
\par\clearpage\null\clearpage
% Unnumbered printed509; principal516, Google520.
\begin{center}
{\Large\textbf{N A C H L A S S.}}\par\bigskip
{\large [Reduktion schiefer Winkel auf den Horizont.]}\par\bigskip
{[}1.]\par\medskip
{\large Reduktion auf den Horizont.}
\end{center}
\[
\begin{array}{ll}
{[}A\ .\ .\ .\ .\ .\ .\ .&\text{schiefer Winkel}\\
A+x\ .\ .\ .\ .&\text{horizontaler Winkel}\\
h,\,h'\ .\ .\ .\ .\ .&\text{Höhen.}
\end{array}
\]
\noindent Man setze]
\[
\begin{gathered}
\tfrac12(A+h+h')=s\\
\tfrac12(h+h')=\sigma,\qquad\tfrac12(h-h')=\delta;
\end{gathered}
\]
\noindent [dann ist]
\[
\begin{aligned}
\operatorname{tang}\tfrac12(A+x)
&=\sqrt{\frac{\sin(\frac12A+\delta)\sin(\frac12A-\delta)}{\cos(\frac12A+\sigma)\cos(\frac12A-\sigma)}}\\
&=\sqrt{\frac{\sin(s-h)\sin(s-h')}{\cos s\cos(A-s)}}.
\end{aligned}
\]
\begin{center}[2.]\par\medskip
{\large Reduktion schiefer Winkel auf den Horizont.}\end{center}
\[
\begin{array}{ll}
A\ .\ .\ .\ .\ .\ .&\text{schiefer Winkel}\\
x\ .\ .\ .\ .\ .\ .&\text{Korrektur}\\
h,\,h'\ .\ .\ .\ .&\text{Höhen}
\end{array}
\]
\[
\frac{\operatorname{tang}\frac12A}{4\,.\,206265}(h+h')^2=M,\qquad
\frac{\operatorname{cotang}\frac12A}{4\,.\,206265}(h-h')^2=N.
\]
\noindent Näherungswert $x=M-N$.

% Printed510; principal517, Google521.
\par\clearpage
\noindent\makebox[0pt][l]{\fontsize{13}{16}\selectfont 510}\hfill{\small\textbf{NACHLASS.}}\hfill\null\par\medskip
\[
\begin{aligned}
\sec\tfrac12x^{\frac13}\,.\,\frac{\sin(A+x)}{\sin(A+\frac12x)}&=\alpha\\[4pt]
\frac{\operatorname{tang}\frac12(A+x)}{\operatorname{tang}\frac12A}&=\beta\\[4pt]
\left(\sec\frac{h+h'}2\right)^{\frac23}&=\mu\\[4pt]
\left(\sec\frac{h-h'}2\right)^{\frac23}&=\nu.
\end{aligned}
\]
\noindent Berichtigter Wert: $\displaystyle x=\frac{\alpha\beta M}{\mu}-\frac{\alpha N}{\beta\nu}$.
\begin{center}Beispiel.\end{center}
\noindent Zwischen Clausberg und Geismar gemessener schiefer Winkel
\[
A=105^\circ27'0\overset{\prime\prime}{,}78,\qquad
\tfrac12A=52^\circ43'30\overset{\prime\prime}{,}39;
\]
\[
\begin{array}{r@{\;}c@{\;}r@{\;}r@{\;}r@{\;}c@{\;}r}
h&=&2^\circ&19'&39\overset{\prime\prime}{,}75&=&8379\overset{\prime\prime}{,}75\\
h'&=&1&4&3{,}17&=&3843{,}17.
\end{array}
\]
\noindent\begin{minipage}[t]{.31\linewidth}\centering
$\begin{array}{r@{\ }r}
h+h'\ .\ .\ .\ .\ .&\underline{4{,}087175}\\
{[}(h+h')^2\ .\ .\ .&8{,}174350\\
\frac{\operatorname{tang}\frac12A}{4\,.\,206265}]\ .\ .\ .&\underline{4{,}202072}\\
M\ .\ .\ .\ .\ .\ .\ .&2{,}376422\\[3pt]
\beta\ .\ .\ .\ .\ .\ .\ .\ .&+477\\
1:\mu\ .\ .\ .\ .\ .\ .&-127\\
\alpha\ .\ .\ .\ .\ .\ .\ .\ .&\underline{-\phantom{1}64}\\
&2{,}376708\\[3pt]
&238\overset{\prime\prime}{,}07
\end{array}$
\end{minipage}\hfill
\begin{minipage}[t]{.31\linewidth}\centering
$\begin{array}{r@{\ }r}
h-h'\ .\ .\ .\ .\ .&\underline{3{,}656728}\\
{[}(h-h')^2\ .\ .\ .&7{,}313456\\
\frac{\operatorname{cotang}\frac12A}{4\,.\,206265}]\ .\ .\ .&\underline{3{,}964958}\\
N\ .\ .\ .\ .\ .\ .\ .&1{,}278414\\[3pt]
1:\beta\ .\ .\ .\ .\ .\ .&-477\\
\makebox[0pt][r]{$-$\qquad}1:\nu\ .\ .\ .\ .\ .\ .&-\phantom{1}17\\
\alpha\ .\ .\ .\ .\ .\ .\ .\ .&\underline{-\phantom{1}64}\\
&1{,}277856\\[3pt]
&18\overset{\prime\prime}{,}96
\end{array}$
\end{minipage}\hfill
\begin{minipage}[t]{.34\linewidth}\centering
$\begin{array}{r@{\ }r}
{[}1:4\,.\,206265]\ .\ .\ .&4{,}083515\\
\operatorname{tang}\tfrac12A\ .\ .\ .\ .\ .\ .\ .&0{,}118557\\
\operatorname{tang}\tfrac12(A+x)\ .\ .\ .&\underline{0{,}119034}\\
{[}\sqrt[3]{\sec\tfrac12x}\ .\ .\ .\ .\ .\ .\ .&0{,}000000\\
\sin(A+x)\ .\ .\ .\ .\ .&9{,}983887\\
\sin(A+\tfrac12x)\ .\ .&\underline{9{,}983951]}\\
M=&237\overset{\prime\prime}{,}91\\
N=&\underline{18{,}99}\\
M-N=&218{,}92
\end{array}$
\end{minipage}
\[
x=+219\overset{\prime\prime}{,}11.\hspace*{.36\linewidth}
\]
\noindent Horizontaler Winkel $=105^\circ30'39\overset{\prime\prime}{,}89$.

% Printed511; principal518, Google522.
\par\clearpage
\noindent\null\hfill{\small\textbf{REDUKTION SCHIEFER WINKEL AUF DEN HORIZONT.}}\hfill\makebox[0pt][r]{\fontsize{13}{16}\selectfont 511}\par\bigskip
\begin{center}[3].\par\medskip
{\large Reduktion eines schiefen Winkels auf den Horizont.}\end{center}
\noindent [Es seien]
\[
\begin{array}{ll}
\alpha\ \text{und}\ \beta\ .\ .\ .\ .\ .\ .&\text{Höhen, in Bogensekunden ausgedrückt}\\
A\ .\ .\ .\ .\ .\ .\ .\ .\ .\ .\ .\ .&\text{[gemessene] Distanz}\\
A'=A+r\ .\ .\ .\ .&\text{Horizontaldistanz.}
\end{array}
\]
\noindent [Dann ist]
\[
2\sin\tfrac12r=\frac{\sin\frac12(\alpha+\beta)^2\operatorname{tang}\frac12A-\sin\frac12(\alpha-\beta)^2\operatorname{cotang}\frac12A}{\cos\alpha\cos\beta}\,.\,\frac{\sin A}{\sin(A+\frac12r)}.
\]
\noindent [Da angenähert
\[
\sin\tfrac12u=\frac{u''}{2\,.\,206265}\sqrt[3]{\cos\tfrac12u}
\]
\noindent gilt, ergibt sich für die Berechnung von $r$ das folgende Schema:]
\[
\begin{array}{r@{\ }l}
\mathrm C\log\cos\alpha=&\\
\mathrm C\log\cos\beta=&\\
\log\mathrm{const}=[\log\frac1{4\,.\,206265}=&\underline{4{,}083\,5149-10]}\\
{[}\log c=]&
\end{array}
\]
\noindent\begin{minipage}[t]{.49\linewidth}\centering
$\begin{array}{r@{\ }l}
\log\operatorname{tang}\tfrac12A=&\\
2\log(\alpha+\beta)=&\\
\tfrac23\log\cos\tfrac12(\alpha+\beta)=&\underline{\hspace{2.7cm}}\\
\log a=&\\
a=&
\end{array}$
\end{minipage}\hfill
\begin{minipage}[t]{.49\linewidth}\centering
$\begin{array}{r@{\ }l}
\log\operatorname{cotang}\tfrac12A=&\\
2\log(\alpha-\beta)=&\\
\tfrac23\log\cos\tfrac12(\alpha-\beta)=&\underline{\hspace{2.2cm}}\\
\log b=&\\
b=&
\end{array}$
\end{minipage}
\par\bigskip
\noindent\begin{minipage}[t]{.34\linewidth}\centering $a-b=$\end{minipage}\hfill
\begin{minipage}[t]{.65\linewidth}\raggedleft
$\begin{array}{r@{\ }l}
\log(a-b)=&\\
{[}\log c=]&\\
\mathrm{Korr.}=\left[\log\left\{\sec\tfrac12r^{\frac13}\,.\,\frac{\sin A}{\sin(A+\frac12r)}\right\}=\right]&\underline{\hspace{1.8cm}}\\
\log r=&\\
r=&
\end{array}$
\end{minipage}
\begin{center}\vspace{1cm}\rule{2.7cm}{.4pt}\end{center}

% Printed512; principal519, Google523.
\par\clearpage
\noindent\makebox[0pt][l]{\fontsize{13}{16}\selectfont 512}\hfill{\small\textbf{BEMERKUNGEN.}}\hfill\null\par\bigskip
\begingroup\small
\begin{center}\vspace{2cm}\textbf{BEMERKUNGEN.}\end{center}
\par\noindent\hspace*{\parindent}Die Notiz [1] findet sich in einem Tagebuch der Sternwarte, II 1813--1817, die Notiz [2] auf der letzten Seite einer Logarithmentafel und die Notiz [3] auf der letzten Seite von \textsc{Gauss}' Exemplar von »G. F. \textsc{Röslers} Handbuch der praktischen Astronomie für Anfänger und Liebhaber, Erster Theil. Tübingen 1788.«
\par\noindent\hspace*{\parindent}Bei [2] lautet im Original der Ausdruck für $\alpha$:
\[
\alpha=\sec x^{\frac13}\,.\,\frac{\sin(A+x)}{\sin(A+\frac12x)},
\]
\noindent während es heißen muss, wie oben angegeben ist. Im Zahlenbeispiel sind einige kleine Rechenfehler verbessert worden.
\par\noindent\hspace*{\parindent}Die der Notiz [2] zugrunde liegende Formel lässt sich wie folgt ableiten. Ist $A$ der schiefe Winkel, $A+x$ der zugehörige Horizontalwinkel, und sind ferner $h$ und $h'$ die Höhenwinkel, also $90^\circ-h$ und $90^\circ-h'$ die Zenitdistanzen, so gilt:
\[
\begin{aligned}
\text{1)}\quad\cos A&=\sin h\sin h'+\cos h\cos h'\cos(A+x)\\
&=\sin h\sin h'\{\cos\tfrac12(A+x)^2+\sin\tfrac12(A+x)^2\}\\
&\qquad+\cos h\cos h'\{\cos\tfrac12(A+x)^2-\sin\tfrac12(A+x)^2\}\\
&=\cos\tfrac12(A+x)^2\cos(h-h')-\sin\tfrac12(A+x)^2\cos(h+h');
\end{aligned}
\]
\noindent subtrahiert man davon
\[
\cos(A+x)=\cos\tfrac12(A+x)^2-\sin\tfrac12(A+x)^2,
\]
\noindent so erhält man die Gleichung
\[
\sin\tfrac12x\sin(A+\tfrac12x)=-\cos\tfrac12(A+x)^2\sin\tfrac12(h-h')^2+\sin\tfrac12(A+x)^2\sin\tfrac12(h+h')^2
\]
\noindent oder
\[
\text{2)}\quad\sin\tfrac12x\,.\,\frac{\sin(A+\frac12x)}{\sin(A+x)}
=\tfrac12\operatorname{tang}\tfrac12(A+x)\sin\tfrac12(h+h')^2-\tfrac12\operatorname{cotang}\tfrac12(A+x)\sin\tfrac12(h-h')^2.
\]
\par\noindent\hspace*{\parindent}In erster Näherung gilt also, wenn $x$, $h$ und $h'$ in Bogensekunden ausgedrückt werden und $\rho=\frac1{206265}$ gesetzt wird:
\[
\begin{aligned}
x&=\tfrac14\rho(h+h')^2\operatorname{tang}\tfrac12A-\tfrac14\rho(h-h')^2\operatorname{cotang}\tfrac12A\\
\text{3)}\qquad&=M-N.
\end{aligned}
\]
\par\noindent\hspace*{\parindent}Da für kleine Winkel näherungsweise
\[
\sin u=\rho\,.\,u''\cos u^{\frac13}
\]
\par\endgroup

% Printed513; principal520, Google524.
\par\clearpage
\noindent\null\hfill{\small\textbf{REDUKTION SCHIEFER WINKEL AUF DEN HORIZONT.}}\hfill\makebox[0pt][r]{\fontsize{13}{16}\selectfont 513}\par\medskip
\begingroup\small
\noindent gilt, folgt aus Gleichung (2) der genauere Näherungswert
\[
\text{4)}\quad x=\left(\sec\frac x2\right)^{\frac13}\,.\,\frac{\sin(A+x)}{\sin(A+\frac12x)}
\GaussOldenburgBracePair{\tfrac14\rho\frac{(h+h')^2}{\left(\sec\frac{h+h'}2\right)^{\frac23}}\operatorname{tang}\tfrac12(A+x)
-\tfrac14\rho\frac{(h-h')^2}{\left(\sec\frac{h-h'}2\right)^{\frac23}}\operatorname{cotang}\tfrac12(A+x)}
\]
\noindent oder mit den Bezeichnungen oben auf S. 510:
\[
x=\alpha\left\{\frac{\beta M}{\mu}-\frac N{\beta\nu}\right\}.
\]
\par\bigskip\noindent\hspace*{\parindent}Die Formel der Notiz [3] ist für die Anwendung noch etwas bequemer, da sie weniger indirektes Rechnen erfordert. Sie lässt sich wie folgt ableiten. Nach Gleichung (1) gilt
\[
\cos(A+x)=\frac{\cos A-\sin h\sin h'}{\cos h\cos h'};
\]
\noindent daher ergibt sich
\[
\begin{aligned}
\cos(A+x)-\cos A=\frac1{\cos h\cos h'}\bigl\{&(\cos\tfrac12A^2-\sin\tfrac12A^2)(1-\cos h\cos h')\\
&-(\cos\tfrac12A^2+\sin\tfrac12A^2)\sin h\sin h'\bigr\}
\end{aligned}
\]
\noindent oder
\[
-\sin(A+\tfrac12x)\sin\tfrac12x=\frac1{\cos h\cos h'}\{\cos\tfrac12A^2\sin\tfrac12(h-h')^2-\sin\tfrac12A^2\sin\tfrac12(h+h')^2\}
\]
\noindent oder
\[
\text{5)}\quad\sin\tfrac12x\frac{\sin(A+\frac12x)}{\sin A}
=\frac1{2\cos h\cos h'}\{\operatorname{tang}\tfrac12A\sin\tfrac12(h+h')^2-\operatorname{cotang}\tfrac12A\sin\tfrac12(h-h')^2\}.
\]
\par\noindent\hspace*{\parindent}Setzt man wieder $\sin u=\rho\,.\,u''\cos u^{\frac13}$, so ergibt sich daraus:
\[
\text{6)}\quad x=\frac{\rho\left(\sec\frac x2\right)^{\frac13}}{4\cos h\cos h'}\,.\,\frac{\sin A}{\sin(A+\frac12x)}
\GaussOldenburgBracePair{\frac{(h+h')^2}{\left(\sec\frac{h+h'}2\right)^{\frac23}}\operatorname{tang}\tfrac12A
-\frac{(h-h')^2}{\left(\sec\frac{h-h'}2\right)^{\frac23}}\operatorname{cotang}\tfrac12A}
\]
\noindent oder, wenn man noch
\[
\text{7)}\qquad\alpha'=\left(\sec\frac x2\right)^{\frac13}\,.\,\frac{\sin A}{\sin(A+\frac12x)}
\]
\noindent und wie vorher
\[
\begin{gathered}
M=\rho\left(\frac{h+h'}2\right)^2\operatorname{tang}\tfrac12A,\qquad N=\rho\left(\frac{h-h'}2\right)^2\operatorname{cotang}\tfrac12A\\
\mu=]\sqrt[3]{\left(\sec\frac{h+h'}2\right)^2},\qquad\nu=\sqrt[3]{\left(\sec\frac{h-h'}2\right)^2}
\end{gathered}
\]
\noindent so erhält man:
\[
\text{8)}\qquad x=\frac{\alpha'}{\cos h\cos h'}\left\{\frac M\mu-\frac N\nu\right\}.
\]
\begin{flushright}\textsc{Krüger.}\end{flushright}
\par\endgroup
\begin{center}\rule{3.5cm}{.4pt}\end{center}
\noindent\makebox[0pt][l]{\small IX.}\makebox[.86\linewidth][r]{65}\par

% Unnumbered printed514; principal521, Google525.
\par\clearpage
\begin{center}
{\Large NACHLASS UND BRIEFWECHSEL.}\par\vspace{1cm}
{\large [Repetitionsbeobachtungen.]}\par\bigskip
{[}1.]
\end{center}
\[
\frac{\text{Wahrscheinlicher Fehler des Ablesens}}{\text{Wahrscheinlicher Fehler des Anzielens}}=\frac1k.
\]
\noindent [$A$, $B$, $C$,\ .\ .\ .\ seien die gleichmäßig verteilten Ablesungen.]
\begin{center}Wahrscheinlichstes Resultat\end{center}
\noindent aus Beobachtungen
\[
\begin{array}{c@{\qquad\qquad}l}
3&\dfrac{(kk+3)(D-A)+(C-B)}{3kk+10}\\[10pt]
4&\dfrac{(kk+2)(E-A)+(D-B)}{4kk+10}\\[10pt]
5&\dfrac{(k^4+5kk+5)(F-A)+(kk+3)(E-B)+(D-C)}{5k^4+28kk+35}\\[10pt]
6&\dfrac{(k^4+4kk+3)(G-A)+(kk+2)(F-B)+(E-C)}{6k^4+28kk+28}.
\end{array}
\]
\par\noindent\hspace*{\parindent}Die Koeffizienten bilden eine rekurrente Folge mit den Rekursionskoeffizienten $kk+2$, $-1$.
\par\bigskip\noindent\hspace*{\parindent}Wird nur zu Beginn sowie nach der ersten, $(n-1)^{\text{ten}}$ und $n^{\text{ten}}$ Beobachtung abgelesen [und sind diese Ablesungen $A$, $B$, $Y$, $Z$], so ist das wahrscheinlichste Resultat
\[
=\frac{((n-2)kk+n)(Z-A)+(n-2)(Y-B)}{n(n-2)kk+nn+(n-2)^2}.
\]
\par\noindent\hspace*{\parindent}Wahrscheinlicher Fehler
\[
=\sqrt{\frac{(n-2)k^4+(2n-2)kk+2}{(nn-2n)k^4+(2nn-4n+4)kk}}.
\]

% Printed515; principal522, Google526.
\par\clearpage
\noindent\null\hfill{\small\textbf{NACHLASS UND BRIEFWECHSEL.\quad REPETITIONSBEOBACHTUNGEN.}}\hfill\makebox[0pt][r]{\fontsize{13}{16}\selectfont 515}\par\medskip
\par\noindent\hspace*{\parindent}Der wahrscheinliche .Fehler, wenn \textls[120]{nur} zu Beginn und am Ende abgelesen wird:
\[
\sqrt{\frac{nkk+2}{nnkk}};
\]
\noindent nach \textsc{Svanberg}s Methode:
\[
\frac{60+6(n(n+2)+2)kk}{5n(n+1)(n+2)kk}.
\]
\begin{center}\vspace{.7cm}[2.]\par\bigskip
{\large\textsc{Gauss} an \textsc{Bessel}.\quad Göttingen, 27. Januar 1819.}\end{center}
\par\noindent\hspace*{\parindent}.\ .\ .\ .\ .\ .\enspace Ich überlege gerade, meiner Theorie der kleinsten Quadrate noch einige Zusätze zu geben. Ein Punkt ist die Behandlung der Beobachtungen mit Repetitionsinstrumenten. \textsc{Laplace} hat darüber kürzlich eine Abhandlung veröffentlicht und unter anderem \textsc{Svanberg}s Verfahren getadelt. Man könnte aber, was \textsc{Laplace} über \textsc{Svanberg} sagt, »c’est un nouvel exemple des illusions auxquelles on est exposé dans ces recherches délicates«, gewissermaßen auf jenen großen Geometer selbst anwenden. Das Verfahren, den ganzen durchlaufenen Bogen lediglich durch $n$ zu teilen (das \textsc{Laplace} in Schutz nimmt), entspricht nur dann der Wahrscheinlichkeitstheorie, wenn die Ablesefehler (und die Teilungsfehler) gegenüber den Anzielfehlern vernachlässigbar sind, besonders wenn die Anzahl der Beobachtungen nicht sehr groß ist. \textls[120]{Jene} Fehler hat \textsc{Laplace} ganz ignoriert. Es ist sehr bemerkenswert, dass umgekehrt, wenn die Anzielfehler gegenüber den Ablese- und Teilungsfehlern vernachlässigbar sind, das der Wahrscheinlichkeitstheorie entsprechende Verfahren ganz mit dem \textsc{Svanberg}schen übereinstimmt. Auch scheint \textsc{Svanberg} nur die letzteren Fehler im Sinn gehabt zu haben (die bei französischen Instrumenten die Anzielfehler wohl auch erheblich überwiegen). Betrachtet man beide Fehler als zugleich vorhanden, so erscheint das Problem sehr schwierig. Geht man es jedoch auf die rechte Art an, ist es äußerst einfach und elementar, wie Sie ohne Zweifel ebenfalls gleich finden werden. Indessen dürfte sich diese Vorgehensweise, wie das Ei des Kolumbus, wohl nicht jedem gleich darbieten, zumal ich allgemein bemerke, dass manche Astronomen, z. B. \textsc{Lindenau} und \textsc{Littrow}, die Methode
\par\medskip\noindent\makebox[.86\linewidth][r]{65\textsuperscript{*}}\par

% Printed516; principal523, Google527.
\par\clearpage
\noindent\makebox[0pt][l]{\fontsize{13}{16}\selectfont 516}\hfill{\small\textbf{BRIEFWECHSEL.\quad REPETITIONSBEOBACHTUNGEN.}}\hfill\null\par\medskip
\noindent der kleinsten Quadrate nicht immer in ihrem wahren Geist anwenden. Aus diesem Grund scheint es mir nicht ganz unverdienstlich, mehrere Punkte dieser Theorie noch eigens auszuarbeiten. Bei 12 Beobachtungen ist, wenn die Teilungsfehler 3-mal so groß sind wie die Anzielfehler, das Gewicht des Resultats
\[
\begin{array}{l@{\;}c@{\;}l}
\text{nach \textsc{Laplace}s Methode}&=&4{,}80\\
\text{nach \textsc{Svanberg}s Methode}&=&7{,}00\\
\text{nach der strengen Methode}&=&7{,}08,
\end{array}
\]
\noindent wobei das Gewicht aus einer Beobachtung \textls[120]{ohne} Ablesefehler $=1$ gesetzt ist.
\par\noindent\hspace*{\parindent}Schon durch Hinzunahme der ersten und elften Ablesung würde das Gewicht nach der richtigen Methode auf 6,26 erhöht.\ .\ .\ .\ .\ .\ .
\begin{center}\rule{2.7cm}{.4pt}\end{center}
\begingroup\small
\begin{center}\textbf{BEMERKUNGEN.}\end{center}
\par\noindent\hspace*{\parindent}Die Notiz über Repetitionsbeobachtungen ist auf dem letzten Blatt von \textsc{Gauss}' Exemplar von \textsc{Laplace}s »Deuxième supplément à la théorie analytique des probabilités, Février 1818« eingetragen. Auf diese Abhandlung bezieht sich auch der unter [2] im Auszug mitgeteilte Brief an \textsc{Bessel}, der nach dem Original abgedruckt wurde. Eine Theorie der Repetitionsbeobachtungen sowohl für gleichmäßig als auch für ungleichmäßig verteilte Ablesungen wurde 1834 von \textsc{Bessel} veröffentlicht (Astronomische Nachrichten, Band XI, Nr. 256, S. 269 ff.).
\par\noindent\hspace*{\parindent}Das Gewicht eines Ablesefehlers sei 1, das Gewicht eines einzelnen Visurfehlers sei $p$. Erfolgen die Ablesungen in regelmäßigen Abständen von je $r$ Repetitionen, so beträgt das Gewicht des Fehlers, der gleich der Summe der Visurfehler zwischen 2 aufeinanderfolgenden Ablesungen ist, mithin $\frac{p}{2r}$. \textsc{Gauss} setzt $\dfrac{p}{2r}=\dfrac1{kk}$.
\par\noindent\hspace*{\parindent}Es bezeichne $x$ den gesuchten Winkelwert, sodass $rx$ aus je 2 aufeinanderfolgenden Ablesungen erhalten wird. Die Ablesungen selbst seien $l_0$, $l_1$,\ldots, $l_n$. Ferner sei $\mu_1$, $\mu_2$, $\mu_3$,\ldots\ eine Folge von Zahlen, die aus der Rekursionsformel
\[
-\mu_{i-1}+g\mu_i-\mu_{i+1}=0,\qquad g=\frac{2r}p+2=kk+2
\]
\noindent unter der Bedingung $\mu_0=0$ und $\mu_1=1$ hergeleitet werden. Es ist also
\[
\begin{array}{r@{\;}c@{\;}l@{\;}c@{\;}l}
\mu_1&=&1&&\\
\mu_2&=&g&=&kk+2\\
\mu_3&=&gg-1&=&k^4+4kk+3\\
\mu_4&=&g^3-2g&=&k^6+6k^4+10kk+4\\
\multicolumn{5}{c}{\text{usw.}}
\end{array}
\]
\par\noindent\hspace*{\parindent}Zur direkten Berechnung der $\mu$ dient die Formel:
\[
k\sqrt{kk+4}\,.\,\mu_i=u^i-u^{-i};\qquad u=\tfrac12(kk+2+k\sqrt{kk+4}).
\]
\par\endgroup

% Printed517; principal524, Google528.
\par\clearpage
\noindent\null\hfill{\small\textbf{BEMERKUNGEN.\quad REPETITIONSBEOBACHTUNGEN.}}\hfill\makebox[0pt][r]{\fontsize{13}{16}\selectfont 517}\par\medskip
\begingroup\small
\par\noindent\hspace*{\parindent}Ist nun die Anzahl der Ablesungen $n=2m$, so ergibt sich allgemein (nach \textsc{Andrae} und \textsc{Zachariae}, Den danske Gradmaaling, Band II, 1872, S. 296 ff.)
\[
x=\frac2p\,.\,\frac{\displaystyle\sum_{\lambda=0}^{m-1}\mu_{m-\lambda}(l_{n-\lambda}-l_\lambda)}{n\mu_{m+1}-(n+2)\mu_m}
\]
\noindent oder auch
\[
rx=\frac{\displaystyle\sum_{\lambda=0}^{m-1}\mu_{m-\lambda}(l_{n-\lambda}-l_\lambda)}{\dfrac1{kk}\{(nkk+n-2)\mu_m-n\mu_{m-1}\}}.
\]
\par\noindent\hspace*{\parindent}Das reziproke Gewicht von $x$ ist:
\[
\frac1P=\frac2{pr}\,.\,\frac{\mu_{m+1}-\mu_m}{n\mu_{m+1}-(n+2)\mu_m}
=\frac{kk}{rr}\left(\frac1n+\frac1n\,.\,\frac{2\mu_{\frac n2}}{(nkk+n-2)\mu_{\frac n2}-n\mu_{\frac n2-1}}\right).
\]
\par\noindent\hspace*{\parindent}Setzt man also
\[
(nkk+n-2)\mu_m-n\mu_{m-1}=kkM,
\]
\noindent so wird
\[
\frac1P=\frac{kk}{rr}\left(\frac1n+\frac{2\mu_m}{nkkM}\right)
=\frac1{rrM}((kk+1)\mu_m-\mu_{m-1})=\frac1{rrM}(\mu_{m+1}-\mu_m).
\]
\par\noindent\hspace*{\parindent}Und ist die Anzahl der Ablesungen $n=2m+1$, so ist
\[
x=\frac2p\,.\,\frac{\displaystyle\sum_{\lambda=0}^{m}(\mu_{m-\lambda+1}+\mu_{m-\lambda})(l_{n-\lambda}-l_\lambda)}{n\mu_{m+2}-2\mu_{m+1}-(n+2)\mu_m}
\]
\noindent oder
\[
rx=\frac{\displaystyle\sum_{\lambda=0}^{m}(\mu_{m-\lambda+1}+\mu_{m-\lambda})(l_{n-\lambda}-l_\lambda)}{\dfrac1{kk}\{(nkk+2n-2)\mu_{m+1}-2(n+1)\mu_m\}}.
\]
\par\noindent\hspace*{\parindent}Hierbei ist das reziproke Gewicht von $x$
\[
\begin{aligned}
\frac1P&=\frac2{pr}\,.\,\frac{\mu_{m+2}-\mu_m}{n\mu_{m+2}-2\mu_{m+1}-(n+2)\mu_m}\\
&=\frac{kk}{rr}\left(\frac1n+\frac1n\,.\,\frac{2(\mu_{\frac{n+1}2}+\mu_{\frac{n-1}2})}{(nkk+2n-2)\mu_{\frac{n+1}2}-2(n+1)\mu_{\frac{n-1}2}}\right).
\end{aligned}
\]
\par\noindent\hspace*{\parindent}Wird mithin
\[
(nkk+2n-2)\mu_{m+1}-2(n+1)\mu_m=kkN
\]
\noindent gesetzt, so ist
\[
\frac1P=\frac{kk}{rr}\left(\frac1n+\frac{2(\mu_{m+1}+\mu_m)}{nkkN}\right)
=\frac1{rrN}((kk+2)\mu_{m+1}-2\mu_m)=\frac1{rrN}(\mu_{m+2}-\mu_m).
\]
\par\noindent\hspace*{\parindent}Schreibt man $A$, $B$, $C$,\ldots\ statt $l_0$, $l_1$, $l_2$,\ldots, so ergeben sich für $n=4$ bzw. 6 und für $n=3$ bzw. 5 die auf S. 514 aufgeführten Resultate, die also für den $r$-fachen Winkel gelten. Die $\mu$ und damit auch die Koeffizienten im Zähler des Ausdrucks für den $r$-fachen Winkel bei geradem $n$ sind die Koeffizienten der Entwicklung des Bruchs $\dfrac1{1-(kk+2)y+yy}$ in eine Reihe nach aufsteigenden Potenzen von $y$. Ist dagegen $n$ ungerade, so sind die Koeffizienten im Zähler des Bruchs für den $r$-fachen Winkel gleich den Koeffizienten der Reihe aus der Entwicklung des Bruchs $\dfrac{1+y}{1-(kk+2)y+yy}$.
\par\endgroup
